\documentclass[10pt,a4paper]{amsart}
\usepackage[margin=19mm]{geometry}
\usepackage{amsmath,amssymb,mathtools}
\usepackage[scr=rsfs,cal=euler]{mathalfa}
\usepackage{xfrac}
\usepackage[unicode,hidelinks,bookmarks=false]{hyperref}
\newcommand{\ie}{i.e.\ }
\newcommand{\cf}{cf.\ }
\newcommand{\resp}{respectively\ }
\newcommand{\Subsection}{\subsection}
\providecommand{\nobreakdash}{\nobreak\hbox{-}}
\setlength{\emergencystretch}{2em}
\multlinegap0pt
\usepackage{amssymb}
\usepackage{mathtools}
\usepackage[shortlabels]{enumitem}
\usepackage{xcolor}
\usepackage{tikz}
\newtheorem{corollary}{Corollary}
\newtheorem{theorem}{Theorem}
\newcommand{\N}{\mathbb{N}}
\newcommand{\Z}{\mathbb{Z}}
\title{Additive systems for $\Z$ are undecidable}
\author{Andrei Zabolotskii}
\date{Source: arXiv:2508.17285v2 (CC BY 4.0)}
\begin{document}
\maketitle

\noindent\emph{Source and licence.} Additive systems for $\Z$ are undecidable. Version arXiv:2508.17285v2 (CC BY 4.0), distributed by arXiv under the Creative Commons Attribution 4.0 International licence, http://creativecommons.org/licenses/by/4.0/, as recorded in the arXiv licence metadata for this exact version and shown on the abstract page https://arxiv.org/abs/2508.17285v2. Full licence notice: \texttt{licenses/arxiv-2508.17285-CC-BY-4.0.txt}.\par\smallskip

\noindent\textit{Editorial scope.} Counted unit: the article's theorem th:Collatz (source0.tex lines 404--407). The article's complete original proof of it is delivered with it (source0.tex lines 408--436, one \texttt{proof} environment of 29 lines). No mathematics is written, repaired or re-derived: the statement and the proof are the article's own lines, in the article's own notation, and no retained line is substituted or re-flowed.  Retained uncounted as the source-local context the proof needs: lines 378--381; lines 397--400.  Nothing was omitted from any retained span, so the package carries no omission record.  No retained line contains a citation, so the wrapper carries no bibliography and no bibliography entry.  No retained line contains a figure environment, an image-inclusion command or drawing code, so no image is delivered and the package declares no asset dependency.  Editorial formatting only, in addition: an AMS \texttt{amsart} preamble that restates the article's own declarations and macros used by the retained lines, namely the environments \texttt{corollary}, \texttt{theorem} on separate counters, together with the standard packages those lines need.  The retained environments print in this document as Corollary 1, Theorem 1 in order of appearance, whereas the article prints the same items with the numbering of its own full document.  The complete native source of the article as distributed in the arXiv e-print for this exact version is bundled unchanged as \texttt{sources/183762/source0.tex}.
\begin{equation}
\label{eq:fn2}
f_i(n) = \lfloor n/b_i \rfloor + f_i(r_{b_i}(n)).
\end{equation}

\begin{corollary}
\label{cor:coll-dyn}
A canonical collection is complete if and only if the trajectory of every integer in the corresponding canonical system eventually becomes constant zero.
\end{corollary}

\begin{theorem}
\label{th:Collatz}
There exists a canonical collection $\mathcal{A}_C$ such that $\mathcal{A}_C$ is complete if and only if the Collatz conjecture is true.
\end{theorem}

\begin{proof}
Let $\mathcal{A}_C$ be the canonical collection with $b_i=4^{i+1}$ for $i\in\N_0$ and $t_i(j) = j-b_if_C(j)$ for $j\in[b_i]\setminus[2]$ and $t_i(1) = 1$. The corresponding canonical system $(f_i)_{i\in\N_0}$ satisfies
\[
f_i(j) = 
\begin{cases}
0 & \text{if } j = 1, \\
j/2 & \text{if $j$ is even}, \\
3j+1 & \text{otherwise,}
\end{cases}
\]
for $j\in[b_i]$.
(Our specific choice of $b_i$ is not the only possible one for our purposes.)

We will now show that $\mathcal{A}_C$ has the property required in the theorem, using Corollary~\ref{cor:coll-dyn}.

Suppose that the $i$th value in a trajectory in this canonical system, $n_i$, falls into the interval $[b_i]$. Then the next value in the trajectory, $n_{i+1} = f_{i}(n_i)$, is either 0 (if $n_i\in[2]$) or $f_C(n_i)$, from the definition of $\mathcal{A}_C$. Therefore $n_{i+1} \leqslant f_C(n_i) \leqslant 3n_i+1 < 3b_i+1 < 4b_i = b_{i+1}$, so $n_{i+1}\in[b_{i+1}]$. By induction, at each further step $i'>i$, the trajectory value $n_{i'}$ will always be in the interval $[b_{i'}]$. Then the trajectory itself will proceed as the Collatz trajectory of $n_i$ unless and until it reaches 1, at which point it becomes constant zero. The Collatz conjecture states that it happens for any positive integer value of $n_i$.

It remains to prove that any number $n\in\Z$ ends up on a Collatz trajectory (by getting into the interval $[b_i]$ on the $i$th step for some $i$) and that every Collatz trajectory is reached from some starting number.

Suppose first that $n$ is positive. Assume it is the $i$th value in a trajectory (possibly $i=0$), i.e.\ $n=n_i$, and $n$ is not already in $[b_i]$, i.e.\ is of the form $n_i=b_i k_i + r_i$, $r_i= r_{b_i}(n_i)$ with $k_i\in\N$. Then, recalling Eq.~\eqref{eq:fn2},
\begin{multline*}
n_{i+1} = f_i(n) = k_i + f_C(r_i) \leqslant k_i + 3r_i+1 = n-r_i - (b_i-1)k_i + 3r_i+1 = \\ = n + 2r_i + 1 + k_i - b_i k_i < n  + 2b_i + k_i - b_i k_i = n + 2 - (k_i-2)(b_i-1).
\end{multline*}
If $k_i>2$ then $n_{i+1}<n$, so $k_{i+1}$ will be less than $k_i$ but will not become negative, so at some step $i'>i$ (repeating the argument for each step) we will get $k_{i'}\leqslant2$. If $k_i=2$ or $k_i=1$ then $n_{i+1}<b_{i+1}=4b_i$, as desired.

Suppose instead that $n$ is negative. We will show that it strictly increases with each iteration until it becomes nonnegative, which is then the case already considered in the previous paragraph. Again, let $n = n_i = b_ik_i+r_i$, $r_i = r_{b_i}(n_i)$ with $k_i$ a negative integer. Then $n_{i+1} = f_i(n) = k_i + f_C(r_i)$. Suppose $k_i<-1$, then $n_{i+1}\geqslant k_i \geqslant b_i(k_i+1) > n$, i.e.\ $n_{i'}$ will increase as a function of $i'\geqslant i$ until $k_{i'}$ becomes nonnegative (and then $n_{i'}$ is nonnegative too, as desired) or equal to $-1$. In the latter case, $n_{i'+1} = f_{i'}(n_{i'})>n_{i'}$ too, because either $n_{i'+1}\geqslant0$ or $r_{i'}\in[2]$, in which case $n_{i'+1}=-1>n_{i'}$.

Finally, the Collatz trajectory of any $n\in\N$ is reached by the canonical system corresponding to $\mathcal{A}_C$ starting from $n\cdot\prod_{i=0}^{\lceil\log_4n\rceil}b_i = n\cdot4^{1+2+\dotsb+(\lceil\log_4n\rceil+1)} = n\cdot2^{(\lceil\log_4n\rceil+1)(\lceil\log_4n\rceil+2)}$.
\end{proof}

\end{document}
